% =====================================================================
%  リソース公開＆共同研究募集（冒頭・アジェンダの直後に配置）
%  ホワイトペーパー / Auto Research (OSS) / 共同研究者募集
% =====================================================================

% ---- ホワイトペーパー公開 ----
\begin{frame}{ホワイトペーパーを公開しました}{「AI-Readyデータ」の考え方を1冊にまとめました}
  \begin{itemize}\small\setlength\itemsep{2.4mm}
    \item きょうの話をさらに深く――\kw{AI-Ready Finance} の考え方を体系的にまとめました。
    \item 2つの壁（\kw{膨大性}・\kw{時系列}）と、\kwg{アウトカムにつなぐ}設計を詳述。
    \item どなたでも\kw{無料}でダウンロードいただけます。ぜひダウンロードしてください。
  \end{itemize}
  \vfill
  % ---- 全幅・一行のダウンロード URL バンド ----
  \begin{tcolorbox}[enhanced,colback=brandnavy!5,colframe=brandnavy,boxrule=0.8pt,arc=1.4mm,
    left=3mm,right=3mm,top=2.2mm,bottom=2.2mm,halign=center]
    {\footnotesize\color{brandgray}ダウンロード}\\[1mm]
    \mbox{{\small\color{brandblue}\bfseries https://katsuyaito.github.io/indx-autoresearch-financeaiready/whitepaper.pdf}}
  \end{tcolorbox}
  \vskip1mm
  \bigstate{「投資情報をAI-Readyにせよ」を、\kwg{もう一段深く}。}
\end{frame}

% ---- Auto Research を OSS 公開 ----
\begin{frame}{Auto Research を OSS で公開}{毎日の自動リサーチを、あなたのラボでも}
  \begin{itemize}\small\setlength\itemsep{2.4mm}
    \item 私たちの\kw{自動リサーチ基盤}をオープンソースで公開しました。
    \item GitHub Actions で\kwg{Web検索→構造化→Issue化→サイト生成}まで自動化。
    \item テーマを差し替えれば、あなたの分野の\kw{毎日の調査}をそのまま回せます。ぜひ使ってください。
  \end{itemize}
  \vfill
  % ---- 全幅・一行のリポジトリ URL バンド ----
  \begin{tcolorbox}[enhanced,colback=brandnavy!5,colframe=brandnavy,boxrule=0.8pt,arc=1.4mm,
    left=3mm,right=3mm,top=2.2mm,bottom=2.2mm,halign=center]
    {\footnotesize\color{brandgray}GitHub リポジトリ}\\[1mm]
    \mbox{{\color{brandblue}\bfseries https://github.com/INDXDev/autoresearch}}
  \end{tcolorbox}
  \vskip1mm
  \bigstate{Star \& Fork 歓迎。\kwg{使って、フィードバック}をください。}
\end{frame}

% ---- 共同研究者の募集 ----
\begin{frame}{共同研究者を募集しています}{2つの仮説を、一緒に検証しませんか}
  \begin{columns}[T]
    \begin{column}{0.52\textwidth}
      \begin{casestudy}{一緒に取り組みたい2つの仮説}
        \footnotesize
        \begin{enumerate}\setlength\itemsep{1.4mm}
          \item \kw{膨大なデータのAI-Ready化}\\
            {\scriptsize\color{brandgray} 壁①：膨大性を読み解き、検索・推論可能にする。}
          \item \kw{時系列データのAI-Ready化}\\
            {\scriptsize\color{brandgray} 壁②：時系列を解釈・検証し、アウトカムにつなぐ。}
        \end{enumerate}
      \end{casestudy}
    \end{column}
    \begin{column}{0.44\textwidth}
      \begin{insight}{こんな方を歓迎}
        \footnotesize
        クオンツ／ML 研究者、金融データに関心のあるエンジニア・学生。\\[1.5mm]
        \kwg{立場・所属は問いません}。
      \end{insight}
    \end{column}
  \end{columns}
  \vskip0.5mm
  \begin{center}
    \begin{tikzpicture}
      \node[fnode blue, text width=12cm, font=\small, align=center]
        {共同研究・コラボのご相談は X（Twitter）\kw{@k1ito} まで DM してください\\[0.5mm]
         {\scriptsize\color{brandmute} x.com/k1ito ── お気軽にどうぞ}};
    \end{tikzpicture}
  \end{center}
\end{frame}
